\begin{titlepage}
\centering
{\Large Notas de entrevista técnica\par}
\vspace{1.0cm}
{\huge\bfseries Entrevista con Luo Fuli (罗福莉): el paradigma Agent, el posentrenamiento y la reestructuración organizacional de los laboratorios de modelos\par}
\vspace{0.6cm}
{\Large Zhang Xiaojun conversa con Luo Fuli\par}
\vspace{0.4cm}
{\large Elaborado a partir del video público del Zhang Xiaojun Podcast, la transcripción hecha en GPU y diagramas conceptuales propios\par}
\vspace{0.3cm}
{\large 2026-04-24\par}
\vspace{0.9cm}
\includegraphics[width=0.88\textwidth,height=0.42\textheight,keepaspectratio]{cover.jpg}\par
\vfill
\begin{tcolorbox}[width=0.92\textwidth, colback=black!2!white, colframe=black!55, sharp corners]
\textbf{Título del video}: 138. Entrevista de 3.5 horas con Luo Fuli: ¡el paradigma de la IA ya cambió radicalmente! OpenClaw, el paradigma Agent depende mucho del posentrenamiento, la asignación de GPU y la igualdad dentro de la organización\par
\textbf{Canal del video}: Zhang Xiaojun Podcast\par
\textbf{Duración del video}: 03:36:36\par
\textbf{Enlace del video}: \href{\videourl}{\nolinkurl{\videourl}}\par
\textbf{Notas sobre el material}: YouTube no ofrece subtítulos; el texto se elaboró a partir de la transcripción con faster-whisper large-v3 en CUDA, los capítulos del video, la portada y figuras didácticas propias. Las tomas fijas de la entrevista no se reutilizan en el texto.\par
\textbf{Página del proyecto}: \href{\repourl}{\nolinkurl{\repourl}}\par
\end{tcolorbox}
\end{titlepage}

\tableofcontents
\newpage
